% Part: incompleteness
% Chapter: incompleteness-provability
% Section: introduction

\documentclass[../../../include/open-logic-section]{subfiles}

\begin{document}

\olfileid[id]{inc}{inp}{int}

\olsection{Pengantar}

Hilbert berpendapat bahwa suatu sistem aksioma bagi struktur matematis,
seperti bilangan asli, tidak memadai kecuali memungkinkan kita menurunkan
semua pernyataan yang benar tentang struktur tersebut. Jika digabungkan
dengan minatnya kemudian terhadap sistem deduksi formal, hal ini menunjukkan
bahwa menurutnya kita harus menjamin bahwa, misalnya, sistem formal yang kita
gunakan untuk menalar tentang bilangan asli bukan hanya konsisten, melainkan
juga \emph{lengkap}, yakni setiap pernyataan dalam bahasanya atau negasinya
!!{derivable}. Teorema ketaklengkapan pertama G\"odel menunjukkan bahwa tidak
ada sistem aksioma semacam itu: tidak ada sistem formal aritmetika yang
lengkap, konsisten, dan !!{axiomatizable}. Bahkan, tidak ada teori matematika
yang ``cukup kuat,'' konsisten, dan !!{axiomatizable} yang lengkap.

Tujuan Hilbert yang lebih penting, yang menjadi inti programnya untuk
membenarkan matematika modern (``klasik''), ialah menemukan pembuktian
konsistensi finitistik bagi sistem formal yang merepresentasikan penalaran
klasik. Dalam kaitannya dengan program Hilbert, teorema ketaklengkapan kedua
G\"odel memberi pukulan yang jauh lebih berat. Seperti teorema
ketaklengkapan pertama, teorema ketaklengkapan kedua dapat dinyatakan secara
garis besar. Secara kasar, teorema itu mengatakan bahwa tidak ada teori
aritmetika yang cukup kuat yang dapat membuktikan konsistensinya sendiri. Kita
harus memahami ``cukup kuat'' sebagai mencakup sedikit lebih banyak
daripada~$\Th{Q}$.

Gagasan di balik pembuktian asli G\"odel atas teorema ketaklengkapan dapat
ditemukan dalam paradoks Epimenides. Epimenides, seorang Kreta, menyatakan
bahwa semua orang Kreta adalah pembohong; bentuk paradoks yang lebih langsung
ialah pernyataan ``kalimat ini salah.'' Pada dasarnya, dengan mengganti
kebenaran dengan !!{derivability}, G\"odel berhasil memformalkan
!!a{sentence} yang, secara tidak langsung, menyatakan bahwa dirinya sendiri
tidak !!{derivable}. Jika !!{sentence} tersebut !!{derivable}, teorinya akan
tak konsisten. G\"odel menunjukkan bahwa negasi !!{sentence} itu juga tidak
!!{derivable} dari sistem aksioma yang dipertimbangkannya. (Untuk bagian
kedua ini, G\"odel harus mengasumsikan bahwa teori~$\Th{T}$ bersifat apa yang
disebut ``$\omega$-konsisten.'' $\omega$-Konsistensi berkaitan dengan
konsistensi, tetapi merupakan sifat yang lebih kuat.\footnote{Artinya, setiap
teori yang $\omega$-konsisten bersifat konsisten, tetapi tidak sebaliknya.}
Beberapa tahun setelah G\"odel, Rosser menunjukkan bahwa asumsi konsistensi
biasa~$\Th{T}$ sudah cukup.)

Tantangan pertama ialah memahami cara membangun !!a{sentence} yang merujuk
pada dirinya sendiri. Untuk setiap !!{formula}~$!A$ dalam bahasa~$\Th{Q}$,
misalkan~$\gn{!A}$ menyatakan numeral yang bersesuaian dengan~$\Gn{!A}$.
Perhatikan artinya: $!A$~adalah !!a{formula} dalam bahasa~$\Th{Q}$,
$\Gn{!A}$~adalah bilangan asli, dan $\gn{!A}$~adalah sebuah \emph{suku}
dalam bahasa~$\Th{Q}$. Jadi, setiap !!{formula}~$!A$ dalam bahasa~$\Th{Q}$
memiliki sebuah \emph{nama}, yakni~$\gn{!A}$, yang merupakan suku dalam
bahasa~$\Th{Q}$; hal ini memberi kita kerangka konseptual tempat
!!{formula}s dalam bahasa~$\Th{Q}$ dapat ``mengatakan'' sesuatu tentang
!!{formula}s lain. Lema berikut dikenal sebagai lema titik tetap.

\begin{lem}
Misalkan $\Th{T}$ adalah sebarang teori yang memperluas~$\Th{Q}$, dan
misalkan $!B(x)$ adalah sebarang !!{formula} yang hanya memiliki variabel
bebas~$x$. Maka terdapat !!a{sentence}~$!A$ sedemikian sehingga $\Th{T}
\Proves!A \liff !B(\gn{!A})$.
\end{lem}

Lema tersebut menyatakan bahwa, untuk setiap sifat $!B(x)$, terdapat
!!a{sentence}~$!A$ yang menyatakan ``$!B(x)$ benar tentang saya,'' dan
$\Th{T}$ ``mengetahui'' hal ini.

Bagaimana kita dapat membangun !!a{sentence} semacam itu? Perhatikan versi
paradoks Epimenides berikut, yang berasal dari Quine:
\begin{quote}
``Menghasilkan kepalsuan bila didahului kutipannya'' menghasilkan kepalsuan
bila didahului kutipannya.
\end{quote}
Kalimat ini tidak merujuk langsung pada dirinya sendiri. Kalimat itu hanya
membuat pernyataan mengenai objek sintaksis di antara tanda kutip, dan dalam
hal ini setara dengan kalimat-kalimat seperti
\begin{enumerate}
\item ``Robert'' adalah nama yang bagus.
\item ``Saya berlari.'' adalah kalimat pendek.
\item ``Memiliki tiga kata'' memiliki tiga kata.
\end{enumerate}
Namun, apa yang terjadi jika frasa ``menghasilkan kepalsuan bila didahului
kutipannya'' didahului oleh versi dirinya sendiri yang dikutip? Kita
memperoleh kalimat semula!{} Singkatnya, kalimat tersebut menyatakan bahwa
dirinya salah.

\end{document}
